% !TeX root = ../guide
\chapter{Figures, Tables, \& Plates}\label{ch:figureandtables}
	\warning{%
		Throughout this Chapter there will be code listings that include lines that show the required packages. 
		These lines should be included in your \LaTeX\ preamble, not in the body of your document.
		To make this easy I have actually included a document specifically to add your packages too. This can be found in the \path{includePackages.tex} file found in the \path{00_LaTeX_Files} folder.}

	\tip{%
		\path{./00_LaTeX_Files/includePackages.tex} includes all the packages required for creating all the examples in this document, some of these will not be necessary for your own thesis, however, I have included comments on all packages for what they are used for. 
		This allows you to make a decision on if you will need them as well. 
		If you are unsure, you may just comment out the line so as to not forget the packages that were originally included.}
	\section{Introduction}
		Figures, Tables, and Plates play a crucial role in conveying complex information effectively within academic documents. 
		This chapter will delve into the intricacies of incorporating these visual elements into your \LaTeX\ document, exploring both basic features and advanced techniques to enhance the visual appeal and clarity of your content.
		While there are many types of figures that one might include in a document, this Chapter specifically focuses on the inclusion of graphic figures (pictures) and the layout of multi-figures (sub-figures).

	\section{Inserting Figures}
		In \LaTeX, figures are included using the \pkg{graphicx}\footfullcite{graphicx} package. 
		The \cmd{includegraphics} command is used to insert an image. 
		Let's consider an example:

		\LTXinputExample{./99_Inclusions/Code/example-figure}

		In this example, the \env{figure} environment is used to contain the image. 
		The optional argument \oopt{htbp} for the \env{figure} environment controls the order in which \LaTeX\ will attempt to insert the float (figure).
		\begin{itemize}
			\item \opt{h} tries to place the float \textbf{here}.
			\item \opt{t} tries to place the float at the \textbf{top} of the page.
			\item \opt{b} tries to place the float at the \textbf{bottom} of the page.
			\item \opt{p} tries to place the float on a float only \textbf{page}.
		\end{itemize}

		The \cmd{centering} command ensures the image is centred horizontally.
		The \opt{width} parameter is used to control the size of the image; in this case, it has been set to \path{0.45\linewidth} which will make it fill a space that is 0.45 times (45\%) the width of the current line.
		The \cmd{caption} and \cmd{label} commands provide a caption and label for referencing, respectively.

		Figures can be formatted to meet specific requirements. 
		The \cmd{subfigure} command from the \pkg{subcaption}\footfullcite{subcaption} package can be used for side-by-side figures:

		\LTXinputExample{./99_Inclusions/Code/example-subfigure}

		This example uses the \env{subfigure} environment to create subfigures within a larger \env{figure} environment. 
		The \cmd{hfill} command adds horizontal space between the subfigures.

		While this section provided a few examples on how to make some figures and subfigures, I would strongly recommend checking out some of the more complex examples of figures shown in \Cref{app:examplefigures}.
		That appendix provides the exact code and examples for creating much more intricate subfigure grids and layouts.
		
	\section{Tables and Tabularx}
		Tables in \LaTeX\ are traditionally created using the \env{tabular} environment.
		However, the \pkg{tabularx}\footfullcite{tabularx} package is particularly useful when you want the table to automatically calculate and adjust its column widths to fit a specific dimension. 
		To make writing tables a little easier, lets define some custom column types in our \path{includeMacros.tex} file:

		\lstinputlisting{./99_Inclusions/Code/tabularx-setup}

		Now, lets implement a table utilizing the \pkg{tabularx} custom columns:

		\LTXinputExample{./99_Inclusions/Code/tabularx-example}

		In this example, the \env{tabularx} environment is used, and the custom column types \opt{C}, \opt{L}, and \opt{R} are applied to the columns. 
		This ensures the content is cleanly centered, left-justified, and right-justified, respectively

	\section{Advanced Table Features}

		To create professional-looking, publication-ready tables, the \pkg{booktabs}\footfullcite{booktabs} package is the industry standard. 
		It provides specialized commands for better vertical spacing and styling of table rules (lines):

		\LTXinputExample{./99_Inclusions/Code/booktabs-example}

		By replacing standard \cmd{hline} commands with \cmd{toprule}, \cmd{midrule}, and \cmd{bottomrule}, you create horizontal rules with appropriate, aesthetically pleasing spacing that adheres to formal typesetting guidelines.

	\section{Additional Packages for Enhanced Table Functionality}

		Depending on your data, several other packages can be employed to heavily enhance table functionality:

		\begin{itemize}
		    \item The \pkg{longtable}\footfullcite{longtable} package allows tables to seamlessly span across multiple pages, which is essential for displaying large datasets.
		    \item The \pkg{multirow}\footfullcite{multirow} package provide targeted commands for merging cells that need to span across multiple rows.
		    \item The \pkg{makecell}\footfullcite{makecell} package enables much more complex table cell layouts, such as forced line breaks within a single cell.
		\end{itemize}
		Each of these packages comes with its own set of commands and options. 
		Let's briefly explore the usage of \pkg{longtable}, \pkg{multirow}, and \pkg{makecell}:

		\begin{lstlisting}[float=ht,caption=,label=lst:additionaltablepackages,style=LaTeXStyle,basicstyle=\small\ttfamily,]
\usepackage{longtable}
\usepackage{multirow}
\usepackage{multicolumn}

% Example Longtable
\begin{longtable}{|c|c|}
	\caption{Longtable Example} \label{tab:longtable} \\
	\hline
	\textbf{Header 1} & \textbf{Header 2} \\
	\hline
	\endfirsthead
	\hline
	\textbf{Header 1} & \textbf{Header 2} \\
	\hline
	\endhead
	Content 1 & Content 2 \\
	Content 3 & Content 4 \\
	\hline
\end{longtable}

% Example Multirow and Multicolumn
\begin{table}[htb]
	\centering
	\begin{tabular}{|c|c|c|}
		\hline
		\multirow{2}{*}{\textbf{Multirow-Col1}} & \multicolumn{2}{c|}{\textbf{Multicolumn-Col2-3}} \\
		\cline{2-3}
		& \textbf{Column 2} & \textbf{Column 3} \\
		\hline
		Content 1 & Content 2 & Content 3 \\
		\hline
	\end{tabular}
	\caption{Example Table with Multirow and Multicolumn}
	\label{tab:multirow_multicolumn}
\end{table}
		\end{lstlisting}

		In these examples, the \env{longtable} environment manages tables that span multiple pages by defining headers and footers that repeat. 
		The \cmd{multirow} command is employed to create cells that span vertically across multiple rows, while \cmd{multicolumn} is used for cells that span horizontally across multiple columns.

	\section{Conclusion}
		This Chapter provided a comprehensive overview of including figures, plates, and tables in your \LaTeX\ document. 
		From the basic insertion of figures to advanced table formatting using powerful packages like \pkg{tabularx}, \pkg{booktabs}, and others, you now have a solid foundation to manage visual data in your thesis.
		While this chapter provided a lot of details on how to add figures that have already been pre-generated, you might want to generate figures on the spot-potentially using data generated directly from other programs (such as MatLab\textsuperscript{\tiny\textregistered}, Python, \textit{etc}.).
		For this advanced topic, \Cref{ch:plotsandgraphs} provides the in-depth workings of the \pkg{pgfplots} package and teaches you how to generate consistent, professional-looking plots and graphs programmatically.